% BHU PYQ -- Professional Exam Template
\documentclass[11pt,a4paper]{article}

\usepackage[a4paper,top=2.5cm,bottom=2.5cm,left=2.8cm,right=2.8cm,headheight=14pt]{geometry}
\usepackage{amsmath,amssymb}
\usepackage[T1]{fontenc}
\usepackage[utf8]{inputenc}
\usepackage{lmodern}
\usepackage{microtype}
\usepackage{fancyhdr}
\usepackage{enumitem}
\usepackage{listings}
\usepackage{hyperref}
\usepackage{lastpage}
\usepackage{parskip}

\hypersetup{colorlinks=true,urlcolor=black,linkcolor=black,
  pdftitle={BVCA-304: Java Applet Programming -- BHU PYQ}}

\pagestyle{fancy}
\fancyhf{}
\renewcommand{\headrulewidth}{0.4pt}
\renewcommand{\footrulewidth}{0.4pt}
\fancyhead[L]{\small\textit{Banaras Hindu University}}
\fancyhead[R]{\small\textit{BVCA-304: Java Applet Programming}}
\fancyfoot[C]{\small Page \thepage\ of \pageref{LastPage}}

\lstset{
  basicstyle=\small\ttfamily,
  frame=single,
  breaklines=true,
  columns=flexible,
  keepspaces=true
}

\newlist{parts}{enumerate}{2}
\setlist[parts,1]{label=(\alph*),leftmargin=*,itemsep=4pt,topsep=3pt}
\setlist[parts,2]{label=(\roman*),leftmargin=*,itemsep=2pt,topsep=2pt}
\newcommand{\pts}[1]{\hfill\textit{[#1]}}

\begin{document}

\begin{center}
  {\large\bfseries BANARAS HINDU UNIVERSITY}\\[2pt]
  {\normalsize B.Voc. Semester III Examination, 2022-23}\\[6pt]
  \rule{0.8\linewidth}{0.5pt}\\[6pt]
  {\large\bfseries Paper: BVCA-304 --- Java Applet Programming}\\[4pt]
  {\normalsize Time: Three hours \hfill Maximum Marks: 70}
\end{center}

\rule{\linewidth}{0.4pt}

\smallskip
\noindent\textit{Instructions: ANSWER ALL SECTIONS, CANDIDATES ARE REQUIRED TO WRITE THEIR ANSWERS IN THEIR}

\bigskip

— REENO057601
COMPUTER APPLICATION
BVCA - 304 : Java Applet Programming
NOTE: ANSWER ALL SECTIONS, CANDIDATES ARE REQUIRED TO WRITE THEIR ANSWERS IN THEIR
OWN WORDS AS FAR AS PRACTICABLE.
Section - A |
Long Answer type questions to be answered in about 750 words each: 15x2=30 Marks

\medskip\noindent\textbf{Question 1.} What are Data types and Variables in Java. Differentiate between static, instance and local variables. Why |
keywords cannot be used as variables?
\centerline{\textbf{OR}}

\medskip\noindent\textbf{Question 2.} Discuss about if else Statement and Switch Statement. Write a program to display days of week fora
number using Switch Statement. (Let us suppose days of week are assigned numbers as Sunday-0, Monday-1, |
Tuesday-2 and so on. So, if user enters 2 value 4, the output should be Thursday)

\medskip\noindent\textbf{Question 3.} a)What are arrays? Explain its types and also discuss the ways to initialize them.
\begin{parts}
  \item Write a program to find the maximum element in an array of 20 integer values. |
OR
\end{parts}

\medskip\noindent\textbf{Question 4.} Explain about for, while and do-while loop structures in Java with example. What is the difference
between entry-controlled loops and exit-controlled loops?

\bigskip\noindent{\large\bfseries SECTION --- B}
\rule{\linewidth}{0.4pt}\smallskip

Short Answer type questions to be answered in about 550 words each Any Six: 5x6=30 Marks

\medskip\noindent\textbf{Question 1.} How to read data using Scanner class in java? State at least three methods defined in the Scanner class to
read data.

\medskip\noindent\textbf{Question 2.} What is concatenation of Strings? Write a program to concatenate two strings.

\medskip\noindent\textbf{Question 3.} Differentiate between break and continue statement with example.

\medskip\noindent\textbf{Question 4.} What are the features of Java Programming Language?

\medskip\noindent\textbf{Question 5.} Write a program to check whether the number is prime or not.

\medskip\noindent\textbf{Question 6.} Define JDK and JVM. Why Java is considered as platform independent language?

\medskip\noindent\textbf{Question 7.} What is a Constructor? Explain its types.

\medskip\noindent\textbf{Question 8.} What is an Object and how is it created? Create a Class Student with at least three attributes and create
two objects of the class Student.

\bigskip\noindent{\large\bfseries SECTION --- C}
\rule{\linewidth}{0.4pt}\smallskip

Answer the following questions in a word: 10x1 Marks i

\medskip\noindent\textbf{Question 1.} 1) Who invented Java Programming?
\begin{parts}
  \item Guido van Rossum
  \item James Gosling
  \item Dennis Ritchie
  \item Bjarne Stroustrup
  PF:
  (2) 5,
  \begin{parts}
    \item \_\_\_ is used to find and fix bugs in the Java programs. | :
  \end{parts}
  \item JVM .
  \item JRE
  \item JDK |
  \item JDB
  \begin{parts}
    \item What does an Eclipse IDE allow? '
  \end{parts}
  \item Editing the program
  \item Compiling the program '
  \item Both Editing and Compiling the program |
  \item Editing, Building, Debugging and Compiling the program.
  \begin{parts}
    \item SDK stands for\_\_
  \end{parts}
  \item Software Design Kit
  \item Serial Development Kit
  \item Software Development Kit
  \item Serial Design Kit
  \begin{parts}
    \item Which of the following is not an OOPS concept in Java?
  \end{parts}
  \item Polymorphism
  \item Inheritance |
  \item Compilation i
  \item Encapsulation
  \begin{parts}
    \item What will be the output of the following Java program? {
    Classvariable\_scope :
    {
  \end{parts}
\begin{lstlisting}[language=C]
publicstaticvoidmain(String ares|]}) |
{
int x; i
x=5; |
{
int y = 6;
System.out.print(x + ““ +y);
}
System.out.printin(x + “” + y); |
}
}
\end{lstlisting}
  \item Compilation error
  \item Runtime error :
  \item 5656
  \item 565
  amt)
  a rr SR
  , @)
  \begin{parts}
    \item What will be the output of the following Java program?
    Class output
    {
  \end{parts}
\begin{lstlisting}[language=C]
publicstaticvoidmain(String args[])
{
StringBuffer s1 = newStringBuffer(“Quiz”);
StringBuffer s2 = s1.reverse(); ‘
System.out.printIn(s2); ;
}
} j
\end{lstlisting}
  \item QuizziuQ
  \item ziuQQuiz
  \item Quiz
  \item ziuQ
  \begin{parts}
    \item In which memory a String is stored, when we create a string using new operator?
  \end{parts}
  \item Stack
  \item String memory
  \item Heap memory
  \item Random storage space
  \begin{parts}
    \item Which keyword is used for accessing the features of a package?
  \end{parts}
  \item package
  \item import
  \item extends
  \item export .
  \begin{parts}
    \item Which of these values is returned by read{} method is end of file (EOF) is encountered?
  \end{parts}
  \item 0 |
  \item 1 : |
  \item -1 i
  \item Null ‘
\end{parts}

\vfill
\rule{\linewidth}{0.4pt}
\begin{center}\small
\href{https://bhu-exam.vercel.app/}{Click here to visit our website}\\[4pt]\href{https://me-aryan.vercel.app/}{Click here to visit our contributor}\\[4pt]\href{https://quashan-me.vercel.app/}{Click here to visit our contributor}
\end{center}

\end{document}
